\cotefile{14}
\scanpage{31}
2. \uline{Familles de classes de cycles \deleted{\illegible{1} \unsure{à équiv.|à équivalence}} \inserted{mod équiv. d'Alb.}}

\uline{Théorème} 2.1. Soient $X^n$ et $Y^m$ comme d'habitude,
et $Z\in\deleted{\unsure{\operatorname{Cyc}}}A^{i}_{\inserted{\mathrm{alb}}}(X\times Y)$, d'où une \deleted{\unsure{hom composé|hom composée|loi compo}} \inserted{application}
$\unsure{\mathfrak{a}|\mathfrak{z}|\mathfrak{t}}\leadsto Z(\mathfrak{a})$ de $\deleted{\unsure{A_0(Y)}}\inserted{A_{0\,\mathrm{alb}}(Y)}$ dans $A^i_{\mathrm{alb}}(X)$. \deleted{Il existe une
application et une seule de $(A_{0\,\mathrm{alb}}(Y)\simeq\mathrm{Alb}^{*}_{\deleted{Y}}(k))$ dans $A^{\unsure{+|i|*}}_{\mathrm{alb}}(X)$
telle que pour tout $\mathfrak{a}\in\mathrm{Alb}_0(Y)$,} or $A_{0\,\mathrm{alb}}(Y)\overset{S_Y}{\simeq}\deleted{\unsure{B}}A_Y^*(k)$,
d'où un hom de groupes
\[
u_Z^*:\underset{\inserted{\substack{\shortparallel\\\mathrm{Alb}}}}{A_Y}^{*}(k)\longrightarrow A^i_{\mathrm{alb}}(X)
\]
\deleted{dont la restriction à $A_Y^{(n)}(k)$ est} défini par la
condition
\[
\deleted{\unsure{S_Y}}Z(\mathfrak{a})=u_Z(S(\mathfrak{a}))
\]
pour tt $\mathfrak{a}$ (resp. pour tt $\mathfrak{a}$ de degré 0,
si on regarde la restriction en degré 0).

\uline{De plus}, si on regarde l'hom \uline{induit}
sur $A_Y(k)=A_Y^0(k)$
\[
u_Z:A_Y(k)\longrightarrow A^i_{\mathrm{alb}}(X),
\]
on a ce qui suit : \deleted{il existe une hom de
variétés de VA
\[
p:A'\longrightarrow A_Y
\]
et un cycle
\[
Z'\in A^i_{\mathrm{alb}}(\unsure{/}
\]
}
(i) Il ne dépend que de la classe de $Z$ dans
$A^{i\,\tau}_{\deleted{\unsure{\mathrm{alg}}}}(X\times Y)$ \margin{[A réserver pour plus tard}

(ii) Il existe une \deleted{\unsure{homo|base}}\inserted{isogénie}
% archive sheet number: 30
